\scanpage{144}
\mspage{18}
ordinale de variables ordinales $\unsure{\varphi|\phi}$, telle que l'on ait
\[
X_\alpha^{\deleted{\unsure{C_0}}}(f)\in\mathrm{Filt}^{\pi}\qquad\text{si}\quad
\pi=\varphi\bigl(\underbrace{\pi(C_0)}_r,\underbrace{\pi(f)}_s,\underbrace{\pi(\alpha)}_t\bigr)
\tag*{(3.6.)}
\]
L'estimation cruciale concerne \unsure{évidemment|êvidemment} le cas $\alpha=1$.
Notons que $X_1^{C_0}(f)\inserted{=X_1(f)}$ est défini comme une somme amalgamée \deleted{\unsure{dans ce qui|dans le qui|de ce qui}} \unsure{(}en général infinie) d'objets $X_D$, pour un ens. $\Delta(C_0,f)=\Delta(f)$ de diagrammes, savoir les diagrammes \unsure{comm,|comm.}
\diagram{
\begin{tikzcd}[ampersand replacement=\&]
\deleted{\unsure{X}}A \arrow[r] \arrow[d,"i"'] \& X \arrow[d,"f"] \\
B \arrow[r] \& \deleted{\unsure{C}}Y
\end{tikzcd}}
avec $i\in C_0$. Pour un \deleted{$A$} $i$ fixé, et un $f$ fixé, \deleted{des} \inserted{le cardinal de l'ens. des} diagrammes de ce type est majoré par celui de $\mathrm{Hom}(A,X)\times\mathrm{Hom}(B,Y)$. \unsure{On a|Or} $A,B\in\mathrm{Filt}^r$, $X,Y\in\mathrm{Filt}^s$, donc il nous faut une majoration de
\[
\mathrm{card}\,\mathrm{Hom}_M(A,X),\qquad A\in\mathrm{Filt}^r,\quad X\in\mathrm{Filt}^s.
\]
\deleted{
\deleted{\unsure{et} \unsure{à|a} \unsure{tout} $\mathrm{Filt}^r$}
\deleted{Or $\mathrm{Filt}^{\pi_0}$ \unsure{engendre} $M$ par $\unsure{\varinjlim|\varprojlim}$ \unsure{filtrantes strictes},}
\inserted{\unsure{on aura|on a}}
\deleted{\unsure{et}}
\[
A=\varinjlim_I A_i,\qquad X=\varinjlim_J X_j
\]
avec $I$ grand devant $\unsure{\pi}$, $J$ grand devant
\[
\mathrm{card}\,I\le r^{\pi_0},\qquad\mathrm{card}\,J\le s^{\pi_0}
\]
(cf loc cit 9.12. \unsure{c|e}). $\unsure{*}$
}
% archive sheet number: 143
